\section{Human review and the human-approved core}
\label{app:recent-human-review}

\subsection{What was reviewed, and why}
We reviewed different subsets of BugBank for different purposes: to find failure modes of the construction, to improve the automated filter, to assess the validity of objections, to check the benchmark's reference lines, and to review selected mathematical follow-ups. None of these rounds is a condition of membership in BugBank. Separately, we selected a human-approved core of 147 dossiers, from which the primary and supplementary localization benchmarks were built. Table~\ref{tab:population-map} lists the studies and their populations; Table~\ref{tab:review-rounds} lists the review rounds.

\begin{table}[htbp]
\centering\small\setlength{\tabcolsep}{4pt}
\begin{tabularx}{\linewidth}{@{}p{.29\linewidth}X@{}}
\toprule
Collection or study & Question it supports\\\midrule
147-dossier human-approved core & What claim, counterexample, formal proof, and source review are available?\\
100 texts in the main evaluation set; 31 additional texts for supplementary evaluation & Can a model select a line marked incorrect in a text?\\
100 proofs; 60 statements & How do instructions change the objections produced and the tokens used?\\
16 Lean repair tasks with the correction given & Can a model complete and prove a specified correction, with or without a supplied counterexample?\\
50 paired tasks per model & How does retrieving earlier objections affect proof development?\\
13 explanation packages & What further mathematics follows from every case in an earlier selected family-and-repair group?\\
\bottomrule
\end{tabularx}
\caption{Each study reports results for its own question and evaluation set. The benchmark measures reference-line selection; the mathematical case studies establish structural theorems and verified repairs.}
\label{tab:population-map}
\end{table}

Nine human reviewers contributed. Every dossier review answered the same three questions:
\begin{enumerate}
    \item Does the objection accurately represent the original source claim and its assumptions?
    \item Does the counterexample satisfy those assumptions and make the claim false?
    \item Does the Lean formalization correctly represent the counterexample?
\end{enumerate}
Two further checks were separate from these questions. Reference-line checks asked whether the lines marked in a benchmark text state the challenged claim. Consequence assessment, which asks what the objection shows about the surrounding argument, was optional; its recorded reviewer and verdict apply only to that question and that version, and an absent assessment is not an approval.

The rounds overlap by record identifier (Section~\ref{app:core-selection}), so their counts must not be added to obtain a number of unique reviewed dossiers.

\begin{table}[htbp]
\centering\small\setlength{\tabcolsep}{3pt}
\begin{tabularx}{\linewidth}{@{}p{.14\linewidth}p{.37\linewidth}p{.21\linewidth}X@{}}
\toprule
Round & Frame and selection & Purpose & Outcome\\\midrule
Initial review (5--12 Sept.\ 2026) & 110 packets from the 401-dossier census, eight reviewers; the assignment rule is not recorded in the accessible records & Failure modes; revised filtering prompt & 100 returns: 61 all YES, 25 with a NO, 14 uncertain\\
Further random 30 & Stratified random sample (seed 20260921) from the 221 of 285 retained records without an earlier review, by seven venues; one reviewer & Check of the filtered material & 27 all YES\\
Targeted 33 & 32 selected mathematical artifacts and the documentary supplement D05; one reviewer & Selected records and research extensions & 30 all YES\\
Fixed30 & Two-stratum probability sample of the 285 retained records: 20 of the 191 primary (seed 20260922) and 10 of the 94 others (seed 20260923); one reviewer & Resource-wide assessment & 29 all YES\\
Mathematical supplements & Three scoped packets (L01, H09, G07) extending earlier dossiers; one reviewer & Scope of the formal extensions & 3 all YES, with consequence\\
Reference-line rounds & 40 drawn uniformly (seed 20260921) from the 279 scorable records with nonempty reference spans; after corrections, a disjoint 40 from the 239 remaining eligible records (seed 3927191731) & Do the marked lines state the challenged claim? & 35 of 40, then 38 of 40\\
Final stratified audit & Stratified random sample without replacement from the 284 retained dossiers: 67 of the 191 primary and 33 of the 93 others; one reviewer; earlier verdicts masked & Population estimate & 91 of 100; weighted 91.0\% (95\% interval 81.0\%--96.5\%)\\
\bottomrule
\end{tabularx}
\caption{Human-review rounds. The dossier reviews answered the three questions above for their own populations; the reference-line rounds asked the separate annotation question described in Section~\ref{app:additional-human-reviews}. Counts are reviews of the listed artifacts, not unique dossiers; reviewer exposure, version evidence, and overlaps are given in the text.}
\label{tab:review-rounds}
\end{table}

\subsection{Initial review: 100 returns from 110 assignments}
\label{app:human-verification}
Eight reviewers returned 100 completed reviews from 110 assignments between 5 and 12 September 2026. All 100 reviewed identifiers belong to the earlier 401-dossier census. One returned workbook covers 30 records in two batches and the other seven cover ten each; the ten assignment workbooks are blank templates, and only completed returns contribute judgments. The rule used to assign the 110 packets is not recorded in the review records available to us, so we do not describe this round as a random sample.

We recompute each outcome from the explicit answers: three YES answers count as approval, any NO as rejection, and CANNOT DETERMINE without a NO as inconclusive. The table preserves the answers as returned; the accompanying comments are kept in the review evidence.

\begin{center}
\begin{tabular}{lr}
\toprule
Initial returned judgment & Records \\
\midrule
All three answers YES & \humanReviewAllYes{} \\
At least one NO & \humanReviewRejected{} \\
Inconclusive, without NO & \humanReviewInconclusive{} \\
\midrule
Completed reviews & \humanReviewReturned{} \\
Assigned without a completed return & \humanReviewMissing{} \\
\bottomrule
\end{tabular}
\end{center}

Of the 76 reviews that approved both the counterexample's mathematics and its Lean translation, 61 also found that the objection accurately represented the original claim and its assumptions; seven did not, and eight were uncertain. The 15 reviews that did not approve the source interpretation reported four omitted hypotheses or domain conditions, one missing definition, one dispute between an asymptotic and a literal reading, one initial-value discrepancy, one wrongly highlighted passage, six sources that were no longer available in full, and one uncertainty about notation or encoding. All 15 dossiers are excluded from the human-approved core.

\paragraph{What the comments contributed.}
The comments identified issues that neither a calculation nor a successful Lean build could settle. For \path{Algorithms_35__sol5}, exchanging the two sides of a bipartition did not create the distinct bipartition the claim required. For \path{Graph_Theory_20__sol4}, the objection omitted a disconnectedness assumption. For \path{Graph_Theory_12__sol5}, the selected Lean theorem refuted a conjunction where the source asserted an implication, and the reviewer pointed to another theorem in the same file that addressed the implication. These observations shaped the later filtering prompt and the review process described below.

Two rows in the batch-01 return have three YES answers but apparently exchanged comments: the comments for \path{mse_3703719__ans3707488} (Answers row 8) describe the singleton matroid-rank example in \path{proofwiki_132447}, and the ProofWiki row (Answers row 4) describes the MSE record's complementary-half-set example. We preserve both returns as written; excluding them leaves 98 reviews with unambiguous comment attribution, including 59 approvals, and both counts are available in the public review tables.

\paragraph{Review conditions and version records.}
The 110 packets included earlier automated assessments in a background appendix. Reviewers were asked to finish their own judgment before reading it, but the returns do not record when that material was read, so we report this round as unblinded. Each completed row supplies one reviewer's judgment and a completion date rather than time spent. The returns did not include hashes of the particular PDF or Lean file read, so applicability to the current version is documented separately (Section~\ref{app:core-selection}). Forty-four members of the human-approved core are supported by this round. The saved ledger preserves workbook hashes, sheet and row references, comments, and the attribution issue.

\paragraph{Derived filter and public evidence.}
The reviewers' observations informed a revised guidance prompt, preserved in \path{bugbank/reviews/filter_prompt/}, for automated filtering across the corpus. The resulting automated audit examined 425 dossiers and retained 285; a subsequent withdrawal left 284. The comments also motivated the comparisons of original and replacement counterexamples in Section~\ref{app:source-admissibility}. Anonymous row-level answers, attribution flags, and cohort information are in \path{bugbank/reviews/original100/answers.json} and \path{bugbank/reviews/cohorts.json}; from the repository root, \texttt{python -B bugbank/reviews/reproduce.py --check} recomputes the recorded counts, including the attribution-clean sensitivity, and the \href{https://github.com/qqliu/CoPEER/blob/main/bugbank/reviews/README.md}{review README} explains the exported evidence.

\subsection{Source assumptions and checked replacements}
\label{app:source-admissibility}
\label{app:matched-admission-pairs}

Like the initial review above (Section~\ref{app:human-verification}), the certificate comparisons in this subsection concern the earlier construction process. The paired sources, witnesses, certificates, and receipts are indexed in \path{bugbank/reviews/source_continuity/PAIR_INDEX.json}.

We compare all 16 records whose counterexamples were replaced during the earlier checks. Five belong to the current 147-dossier core, including the sorting and Gaussian-integer examples below. All 16 pairs retain the same stored problem and source context, and 15 also retain the same challenged quote, so the change in the counterexample is directly inspectable.

The original Lean files were recovered for all 16 records. Four have saved original build and axiom-check receipts; the other 12 were recompiled unchanged under Lean 4.30.0 and 4.34.0, and all 24 compilations succeeded. All 16 replacement certificates also have build receipts. The replays establish that the recovered files compile; the original receipts document which checks were performed at the time. Three pairs preserve the same quote and have both original and replacement build receipts: the endpoint-order example below, an equality interpreted as an eventual asymptotic claim, and a covariance claim tested with an ensemble different from the source's. The fourth pair with original receipts changes its challenged quote. The public pair index records these distinctions alongside the source excerpts, witnesses and audit reasons.

An inspection of all 123 theorem and lemma declarations in the 12 recompiled files found 78 using no axioms, 42 using only \texttt{propext}, and three using \texttt{propext} and \texttt{Quot.sound}; none uses \texttt{sorryAx} or native reduction. Three inspection drivers also queried comment text as declaration names, and an independent review checked that their outputs still covered every genuine declaration. The files, commands and complete outputs are in \path{bugbank/reviews/source_continuity/earlier_versions/}.

The replacements address different issues. Four original modules already contained an alternative that satisfies the source assumptions: the arXiv probability family, the ProofWiki induction instance, the six-team tournament and the limit calculation at $x=3/2$; the first two already select that stronger formal statement, and the other two contain it alongside the disputed statement. Other replacements restore omitted assumptions or conventions, as in the two examples below. Their interpretation follows the recorded audits, and their membership in this comparison is determined by the recorded replacement rather than by random sampling.

\paragraph{An executable check of the omitted source condition.}
The gap-sorting source explicitly assumes $x_0<x_n$. Its challenged local step says that after the first exchange with $x_m$, the final gap-$m$ sorted array begins with that entry $x_m$. The original witness uses $m=2$ and $A=(3,4,2,5,1)$. Its Lean-checked theorem correctly shows that the final first entry is $1$ rather than $x_2=2$. But $3<1$ is false: this array violates the source's endpoint hypothesis. The replacement $A'=(3,5,2,6,1,7,4)$ satisfies $3<4$; its first gap-2 insertion exchanges $3$ and $2$, and the final array is $(1,5,2,6,3,7,4)$, again beginning with $1$ rather than $x_2=2$. The selected replacement theorem proves the endpoint inequality, the exchange hypothesis and the failed conclusion together.

A supplementary Lean file retains the current certificate's exact bytes as its prefix and adds three axiom-free lemmas. For the explicit predicate $H(a):=a_0<a_{|a|-1}$, they prove $\neg H(A)$, rule out the conjunction of $H(A)$ with the encoded objection conditions, and prove that $A'$ satisfies those conditions. The file compiles under Lean 4.30.0 with exit code zero. This checks the role of the particular source condition; the correspondence between the quoted prose and $H$ remains an inspectable interpretation, and the file does not certify the entire source proof or claim that its final theorem is false.

\paragraph{An equality case that satisfies the source assumptions.}
In \path{mse_605342__ans609345}, the Gaussian-integer source first chooses $\alpha$ from the lattice generated by $1+i$ with $N(a-\alpha b)\leq N(b)$, then claims that coprimality makes this strict. The old certificate takes $a=b=-1$, $\alpha=i$, and correctly computes $2>1$, outside that weak-bound premise. The current Mathlib certificate keeps $a=b=-1$ and uses $\alpha=1+i$. Its theorem explicitly establishes coprimality, opposite coordinate parities, the allowed multiplier and the weak bound, but proves
\[
  N(a-\alpha b)=N(b)=1,\qquad \neg\bigl(N(a-\alpha b)<N(b)\bigr).
\]
A companion theorem proves $b\mid a$ together with coprimality: divisibility by this unit is no contradiction. The repaired objection therefore identifies the equality boundary in the local strict-descent argument; it does not establish nontermination of the algorithm.

Together, these comparisons show the source check that Co-PEER applies: state the relevant assumptions, prove that the counterexample satisfies them, and establish the contradiction. The human judgments and the certificate comparisons are complementary evidence for these three steps.


\subsection{Further random, targeted, and resource-wide reviews}
\label{app:additional-human-reviews}

\paragraph{Further random 30.}
After automated filtering with the revised prompt, one reviewer, who used AI tools, evaluated a stratified random sample of 30 dossiers from the 221 of 285 retained records with no earlier human review. With seed 20260921, the sampling script drew from sorted record identifiers within seven source venues and then shuffled the review order. The sample comprised four Math StackExchange records, five ProofWiki, four QEDBench, seven Reddit, three reference-proof additions, three Stack Exchange, and four arXiv records. The reviewer approved 27 dossiers on all three questions; the individual judgments were 28 YES and 2 NO for source connection (R012, R026), 30 YES for valid mathematics, and 29 YES and 1 NO for Lean statement match (R019). The record notes that the AI assistant saw the prior rationale on its first reading of R001, so we do not claim that earlier judgments were masked in this round. Because the sampling fractions differ across venues, the unweighted approval count describes the reviewed sample rather than a design-weighted estimate for all retained dossiers.

\paragraph{Targeted 33.}
One reviewer returned judgments for 33 explicitly selected artifacts: 32 mathematical artifacts and the documentary supplement D05. Thirty received YES on all three questions, 29 mathematical artifacts and D05; the three exceptions were E01 (source connection NO), E02 (mathematics NO), and S01 (mathematics not applicable). Fourteen of the artifacts correspond to current record identifiers; the others are research and scope reviews, so the artifact count is not a count of distinct dossiers.

\paragraph{Fixed30.}
A separate round evaluated a two-stratum probability sample of the then-current 285-record roster. It drew 20 records from the 191-record primary stratum with seed 20260922, retained unchanged, and 10 from the 94-record complement with seed 20260923, both without replacement from sorted record identifiers. One reviewer used AI tools and personally checked the judgments, reporting prior project exposure for six of the 30 records. All 20 primary records and nine of the ten complement records received YES on all three questions; for I025, the objection did not accurately represent the source claim. The prespecified population weights were 191/285 and 94/285, and 29 of 30 describes the reviewed sample. This frame differs from the later 284-record audit frame, whose complement contains 93 records.

\paragraph{Mathematical supplements.}
Three scoped supplements (L01, H09, G07) received YES on the three questions and on the optional consequence from one reviewer. L01 extends the density analysis of the layered projective-cube construction to all dimensions, with the cited Long--Wagner cube-freeness theorem as an explicit premise; H09 characterizes when the proposed handshaking induction step works while retaining the truth of the handshaking lemma; G07 proves uniqueness of the polynomial representing the fixed fake-Gaussian quotient and its negative degree-43 coefficient, without a minimality claim. They extend earlier mathematical dossiers rather than adding independent records, and the returned packet hashes match the assigned supplements.

\paragraph{Reference-line rounds.}
These checks asked whether the marked lines of a benchmark text state the challenged claim, whether any such line was missing, and whether an unrelated line was marked; the acceptance criterion was at least 38 of 40 complete line selections with answers YES, NO, NO. The first round sampled 40 records uniformly without replacement (seed 20260921) from the 279 records of the 285-record roster that were eligible for scoring and had nonempty reference spans, and approved 35 of 40. After corrections, a fresh sample of 40 was drawn from the 239 remaining eligible records (seed 3927191731), disjoint from the first; it approved 38 of 40 and met the criterion. The two remaining concerns were an omitted line (G14) and an extraneous marked line (G12). These judgments concern the marked lines in the reviewed source versions and do not themselves validate the mathematical objections. Both sets of judgments are preserved in \path{bugbank/reviews/reference_lines/initial40.json} and \path{bugbank/reviews/reference_lines/fresh40.json}.

Records E01, E02, R012, R019, R026, and I025 are all excluded from the human-approved core.

\subsection{Final random audit with masked prior judgments}
Finally, we drew a stratified random sample without replacement from the 284 retained dossiers: 67 of the 191 primary records and 33 of the 93 other records. One reviewer used AI tools and personally checked the judgments on accurate source interpretation, witness validity, and correspondence with the Lean statement. The review packets displayed the source and certificate text and masked earlier verdicts and interpretive commentary; the reviewer reported prior project exposure for 28 records.

All three criteria received a YES for 63 of 67 primary and 28 of 33 other records, 91 of 100 in total. Weighting each stratum by its share of the 284-dossier collection gives an estimated approval rate of 91.0\%, with a simultaneous 95\% sampling interval of [81.0\%, 96.5\%], obtained by combining two exact 97.5\% finite-population intervals by Bonferroni. The component approval counts were 95 of 100 for accurate source interpretation, 98 for witness validity, and 94 for Lean correspondence; the optional reference-line check received 98 YES and two not-applicable responses. The interval summarizes sampling uncertainty for this review procedure on the 284-record population and does not account for reviewer error. The sampling plan and interval method were specified before review; the reviewer-admissibility clarification and the weighted point estimate followed the returned judgments. The design, chronology, and analysis are preserved in \path{bugbank/reviews/random100/PROTOCOL.md}, \path{bugbank/reviews/random100/design.json}, and \path{bugbank/reviews/random100/saved_analysis.json}.

\subsection{Overlaps, final dispositions, and the human-approved core}
\label{app:core-selection}
The rounds overlap by record identifier, not necessarily by identical reviewed versions. The further random 30 share no identifiers with the initial returns; Fixed30 shares seven with the initial returns and two with the further random 30; the final audit shares 21 with the initial returns, 12 with the further random 30, and seven with Fixed30. The 14 targeted artifacts with current record identifiers overlap the initial returns, the further random 30, Fixed30, and the final audit in 5, 2, 0, and 7 records.

All nine nonapprovals in the final audit are excluded from the core: four from its primary stratum (A001, A002, A023, A043) and five from its other stratum (A072, A075, A085, A093, A099). Table~\ref{tab:final-human-audit-exclusions} preserves their three answers and recorded follow-up categories; a previously recorded concern retains its original NO judgment.

\begin{table}[ht]
\centering\small
\setlength{\tabcolsep}{3pt}
\caption{The nine nonapprovals in the final random audit. Q1, Q2 and Q3 check the source interpretation, witness and Lean statement. All nine are excluded from the current 147 dossiers and both benchmark subsets. ``Earlier set'' refers to the original 191 primary and 255 scorable inputs. Certificate versions are identified by SHA-256 prefixes; the linked version ledger gives full source, packet and certificate hashes. Concerns summarize the recorded reviews, not new adjudications.}
\label{tab:final-human-audit-exclusions}
\begin{tabular}{lll p{5.1cm} l}
\toprule
Packet & Earlier set & \shortstack{Negative\\checks} & Recorded concern & \shortstack{Certificate\\version} \\
\midrule
A001 & Primary & Q3 & Connection from the source's group construction to Lean is left to prose. & \texttt{51284638c1b1} \\
A002 & Primary & Q2, Q3 & Zero witness conflicts with an implicit nonzero assumption. & \texttt{af06bbfe5425} \\
A023 & Primary & Q1, Q2, Q3 & Category-label mismatch is treated as a refutation under a disputed translation. & \texttt{c8e308ad5c7d} \\
A043 & Primary & Q3 & Lean supplies the size-one transition law rather than deriving it. & \texttt{76b30ad2964d} \\
A072 & Other scored & Q3 & Boolean instance lacks formal syntax and substitution linking it to the source. & \texttt{d47105809371} \\
A075 & Other scored & Q1, Q3 & Humorous runtime claim is unspecified; Lean checks a supplied runtime function. & \texttt{af1c203b5790} \\
A085 & Other scored & Q1 & True equalities are displayed; the false step is implied rather than asserted. & \texttt{b0893e69a1ae} \\
A093 & Not scored & Q1 & Target is a typing requirement in the model setup, rather than an assertion. & \texttt{0e7eea1a1635} \\
A099 & Not scored & Q1 & Quoted supposition introduces a case; the gap is its later generalization. & \texttt{cf26d14808fe} \\
\bottomrule
\end{tabular}
\end{table}

The version ledger at \path{bugbank/reviews/version_resolution/README.md} links every case to its full artifact hashes, source text, selected Lean declaration, original review and final disposition. For all nine, the certificate bytes and source-line arrays match the reviewed versions; exclusion does not overturn the negative judgments. Seven scored cases received YES on the separate line check, which did not resolve their source, witness or Lean concerns.

To enter the human-approved core, a dossier must have a completed review that answered YES to all three questions and applies to the current version. It must also have no retained negative or uncertain human judgment, and no unresolved concern about whether the objection applies to the selected source claim under its assumptions. Among the 162 dossiers with at least one all-YES review, 13 also had an adverse human judgment (ten from the initial returns, two from the targeted review, and one from the final audit), and two had such unresolved concerns. Excluding these 15 leaves the 147-dossier core, drawn from 137 distinct sources. Of the 147, 113 have exact matches between the reviewed scientific contents and the current contents; for the other 34, we confirmed that the saved reviews apply to the current versions, a confirmation recorded separately from exact content matching. The membership ledger and continuity receipt are \path{bugbank/core/membership.json} and \path{bugbank/core/author_continuity_receipt.json}. Across the earlier 284-dossier roster, 31 dossiers have an adverse human review; none enters the human-approved core. Within the core, 103 dossiers meet the primary-input criteria and 134 are scorable before the separate reference-line exclusions below.

The 100 primary benchmark inputs include 69 with a positive check of the current reference-line selection: 55 have a positive check with the exact current source text and labels, and 27 have a separate positive line-selection review, with overlap between these groups. The older line-review exports identify the reviewed labels but do not establish complete source-version identity. Three dossiers with adverse judgments on the current line selections remain in the human-approved core and are excluded from both evaluation subsets: \path{arxiv_2207.06033v2__source_enriched_codex}, \path{proofwiki_12198__codex_expansion}, and \path{reddit_mathematics_1dz93wr}. This leaves 100 primary and 131 total scorable inputs. The reference-line evidence is preserved in \path{benchmark/core/reference_line_coverage.json}.

All review-round answers are available under \path{bugbank/reviews/original100/}, \path{bugbank/reviews/earlier_selected30/}, \path{bugbank/reviews/selected33/}, \path{bugbank/reviews/selected30/}, \path{bugbank/reviews/scoped_supplements/}, and \path{bugbank/reviews/random100/}; \path{bugbank/reviews/README.md} explains the files and reproduction commands. Section~\ref{app:verified-census} describes the core's documentation.
